\begin{table}[t]
\centering
\caption{Source corpus families, binary label inheritance, and split contributions in the final v1 dataset. Counts are taken from \texttt{data/v1/manifest.json}.}
\label{tab:dataset_sources}
\footnotesize
\setlength{\tabcolsep}{4pt}
\renewcommand{\arraystretch}{1.15}
\begin{tabularx}{\textwidth}{>{\raggedright\arraybackslash}p{0.16\textwidth}>{\raggedright\arraybackslash}p{0.18\textwidth}>{\raggedright\arraybackslash}p{0.28\textwidth}X}
\toprule
Corpus family & Role in thesis & Binary label mapping & Split contribution \\
\midrule
BIPIA (\texttt{email}, \texttt{table}) & Indirect prompt-injection family & Attack builder rows map to $y=1$ as indirect prompt injection. Clean-context rows map to $y=0$ and are retained as benign. & train: 173,439\newline(172,500 attack / 939 benign)\newline test\_main: 29,175\newline(29,025 attack / 150 benign) \\
JailbreakDB family & Direct jailbreak + benign hard negatives & Source \texttt{jailbreak}/\texttt{label} fields are normalized to the binary task. Existing tactic metadata is preserved in \texttt{attack\_type}; otherwise attacks default to \texttt{jailbreak}. & train: 27,837\newline(13,917 attack / 13,920 benign)\newline val: 6,106\newline(3,024 attack / 3,082 benign)\newline test\_main: 6,055\newline(3,058 attack / 2,997 benign) \\
JailbreakBench behaviors & Held-out direct-jailbreak stress split & Harmful behaviors map to $y=1$ as direct jailbreak; benign behaviors map to $y=0$. This family is kept separate as \texttt{test\_jbb}. & test\_jbb: 197\newline(99 attack / 98 benign) \\
\bottomrule
\end{tabularx}
\par\vspace{0.3em}
{\footnotesize Counts are reported as total examples, followed by attack / benign counts in parentheses.}
\setlength{\tabcolsep}{6pt}
\renewcommand{\arraystretch}{1}
\normalsize
\end{table}
